\documentclass{article}
\usepackage{cancel}
\usepackage{gensymb}
\usepackage{tikz}
\usepackage{amsmath}
\usepackage{ amssymb }
\usepackage{listings}
\usetikzlibrary{angles,quotes}
\usetikzlibrary{arrows.meta,calc}


\renewcommand{\thesubsubsection}{\alph{subsubsection})}

\def\hwNUM{10}


\title{Astro HW \hwNUM}
\author{Pierson Lipschultz}

\begin{document}
\maketitle
\setcounter{section}{\hwNUM}

{\large \textbf{NOTE: Work on the page itself is kinda sparse, but all the code used for the exact results is at the end}\footnote{pls dont dock me for not showing enough work ;-;}}
\subsection{Supernova! Wahoo! Yipee!}

\subsubsection{Tme since SN}
Given \dots
\[\theta = \frac{D}{d}\]
We find \dots
\[D_{snr} = (\theta_{snr} \times d_{snr}), \qquad R_{snr} = \frac{D_{snr}}{2} \]
\[t_0 = R_{snr}/V_{snr}\]
\[
\boxed{
    t_0 \approx 3.269 \times 10^9 \text{s} \approx 1.03 \times 10^2 \text{yr}
}
\]

\subsubsection{Neutrinos per sec}
Given\dots
\[F = \frac{L}{4\pi d^2}\]
We find that
\[F = 1.156 \times 10^{16} \text{m}^2\]
\[F_t = \frac{F}{\Delta T}\] 
\[\boxed{F_t \approx 7.712 \times 10^{14} \text{ Neutrinos/s}^2}\]

\subsubsection{Magnitudes and extinction}
Given \dots
\[M_{bol} = M_{bol\odot} - 2.5 \log(L_{bol}/L_\odot),\qquad M_{bol\odot} = 4.74 \qquad L_{bol} = 10^9 L_\odot\]
\[m = M + 5\log(d) - 5, \qquad d = 8.5\text{kpc} \]
\[m_{obsc} = M + 5\log(d) - 5 + A_v\ \qquad A_v = 30 \]
We find \dots
\[M_{bol} = 4.474 -2.5\log(10^9 L_\odot) \approx -17.76\]
\[m_{bol} = -18.1 + 5 \log(8500\text{pc}) -5 \approx -3.113\]
\[m_{Obscbol} = -18.1 + 5 \log(8500\text{pc}) -5 + 30 \approx 26.886\]
\[\boxed{M_{bol} \approx -17.76, \qquad m_{Obscbol} \approx 26.887}\]
This means it is completely impossible to resolve with the naked eye, and, while resolvable with some scopes, still quite difficult.

\subsection{Stellar remnants and angular momentum}
Given \dots
\[I = \frac{2}{5} MR^2\]
\[L = I\omega\]
\[\omega = \frac{\theta}{t}\]

\[P_{sun} = 25.4 \text{days}, \qquad r_{wd} = 6000 \text{km}\]
We find \dots
\[I_{sun} = 2/5 m_\odot r_\odot^2\]
\[\omega_{sun} = (2 \pi)/P_{sun} \]
\[L_{sun} = I_{sun}\omega_{sun} \approx  1.1\times 10^{42}\text{m}^2\text{kg}/\text{s}\]

\[I_{wdSun} = 2/5 m_\odot r_{wd}^2\]
\[\text{By conservation of angular momentum\dots}\]
\[\omega_{wdSun} = L/I_{wdSun}\]
\[P_{wdSun} = (2 \pi / \omega_{wdSun})\]
\[\boxed{P_{wdSun} \approx 163.2\text{s}}\]

\subsection{The nature of pulsars}
\subsubsection{Orbiting WDs}
Given\dots
\[P^2 = \frac{4\pi^2}{G(m_1 + m_2)}\]
\[M_{wd} = .7M_\odot, \qquad r_{wd} = 7 \times 10^6 \text{m}, \qquad P_{wd} = 1.337 \text{s}\]

We find that \dots
\[a = \left(\frac{(p_{wd}^2)(2GM_{wd})}{r\pi^2}\right)^{1/3}\]
\[\boxed{a \approx 2\times 10^6 \text{m}}\]
\subsubsection{Is this possible?}
No, this is not possible. This is because \(\boxed{a < r_{wd}}\), meaning that the WDs semimajor axis would be literally within each other, and thus they would have to be merged. 

\subsubsection{WD Period}
Given \dots
\[v_{rot}= \frac{2\pi R}{P} \rightarrow P = \frac{2\pi R}{v_{rot}}\]
\[v_{rot}(P,M) \approx .1c\left(\frac{P}{1s}\right)^{-1}\left(\frac{M}{.7M_\odot}\right)^{-1/3} \]
We find that \dots

\[v_{rot} \approx 2.662 \times 10^7 \text{m/s}\]
\[P = 2\pi \frac{r_{wd}}{v_{rot}}\]
\[\boxed{P \approx 1.652\text{s}}\]

\subsubsection{Breakup Speed}
Given \dots
\[a_c = \frac{v^2}{r}\]
\[g = \frac{GM}{R^2}\]
We find \dots
\[v_{crit} \rightarrow \frac{v^2}{r} = \frac{GM}{R^2}\]
\[\frac{v^2}{\cancel{r}} = \frac{GM}{R^{\cancel{2}}}\]

\[\boxed{v = \sqrt{\frac{GM}{R}}}\]
\subsubsection{Breakup speed from alternative function}
Given\dots
\[v_{crit} = \sqrt{\frac{GM_{wd}}{r_{wd}}}\]
\[M_{wd} = .7M_\odot, \qquad r_{wd} = 7 \times 10^6 \text{m}\]
We find that \dots
\[\boxed{v_{crit} \approx 3.643 \times 10^{6} \text{m/s}}, \qquad \boxed{v_{crit} \approx 3.583 \times 10^{6} \text{m/s}}\]

\subsubsection{Is this feasible?}
No shot\dots This is because \(v_{rot} > v_{crit}\), meaning that a WD rotating that fast would break up. 

\subsubsection{What if it was a NS?}
Given \dots

\[P = 2\pi \frac{r_{NS}}{v_{rot}}\]
\[v_{crit} = \sqrt{\frac{GM_{NS}}{r_{NS}}}\]
\[M_{ns} = 1.4 M_\odot, \qquad r_{NS} = 10\text{km}\]
We find\dots
\[\boxed{v_{rot} \approx 4.7 \times 10^{4}\text{m/s} }\]
\[\boxed{v_{crit} \approx 1.36 \times 10^{8} \text{m/s}}\]

This means that a NS spinning with that rate \textit{does} make sense, this is because \(\boxed{ v_{rot} < v_{crit}}\), meaning that the NS will \textit{not} break apart, and could have a pulsation period of 1.337s.

\subsection{Mass models of the Milky Way Galaxy}
Given \dots

\[v**2 = \frac{GM}{R} \rightarrow m = \frac{v^2 r}{G}\]

We find \dots
\[m_0 = \frac{v_0^2 r_0}{G}, \qquad m_1 = \frac{v_1^2 r_1}{G}, \qquad m_2 = \frac{v_2^2 r_2}{G}, \qquad m_3 = \frac{v_3^2 r_3}{G}\]

\begin{center}
    \begin{tabular}{|c|c|c|}
        \hline
        \(v_c\)(km\(^{-1}\)) &\(r\)(kpc)&\(M(r)(kg)\) \\
        \hline
        255 & 0.4 & \(1.20\times 10^{40}\) \\
        192 & 2.8 & \(4.77\times 10^{40}\) \\
        220 & 8   & \(1.79 \times 10^{41}\) \\
        235 & 17  & \(4.34\times 10^{41}\) \\
        \hline
    \end{tabular}
\end{center}

\subsection{Galactic center black hole and S0-2}
\subsubsection{}
Given\dots
\[M_{bh} = \frac{4\pi^2a^3}{Gp^2}\] 
\[a = 960\text{AU}, \qquad p = 15.6\text{yr}\]
We find \dots
\[M_{bh} \approx 7.22914 \times 10^{36} \text{kg}\]
\subsubsection{}
Given\dots
\[r = a(1\pm e)\]
\[v_a = \left(\frac{GM_{bh}}{a} \frac{1-e}{1+e}\right)^{1/2}\]
\[v_p = \left(\frac{GM_{bh}}{a} \frac{1+e}{1-e}\right)^{1/2}\]
\[a = 960 \text{AU}, \qquad M_{bh} = 7.22914 \times 10^{36} \text{kg}, \qquad e =  .867\]
We find\dots
\[\boxed{r_a \approx 2.618 \times 10^{14} \text{m}}, \qquad \boxed{r_p \approx 1.91 \times 10^{13}\text{m}}\]
\[\boxed{v_a \approx 4.892 \times 10^{5}\text{m/s} \approx .16\% \text{ speed of light}}\]
\[\boxed{v_pe \approx 6.867 \times 10^{6}\text{m/s} \approx 2.29\% \text{ speed of light}}\]

\subsubsection{}
Given \dots
\[d = r\left(\frac{2M}{m}\right)^{1/3}\]
\[m = 17.5 M_\odot, \qquad r = 7.4R_\odot, \qquad M_{bh} = 7.22 \times 10^{36} \text{kg}\]
We find \dots
\[\boxed{d_{rl} \approx 3.842 \times 10^{11}\text{m}}\]
\[\frac{r_p}{d} \approx 50\]
This star is currently \textit{not} in danger of TDe, this is because \(\boxed{d_{pe} > d_{rl}}\), meaning it doesn't come within the BHs RL, and thus is not at risk.




\end{document}